\section{Contextual response measurement: unfinished development extension}
\label{app:contextual}

This October extension examines the measurement needed to connect emoji
information with ordinary response probabilities. It uses no CLDR targets,
readout-state replacement, or internal intervention. The materials are 24
author-written scenarios in six recurring families: temperature, setting,
completion, direction, forecast, and affect. Every development stage reuses
these scenarios. Repeated conditions increase measurement coverage, not the
number of independent scenarios. Analyses are descriptive, with no population
p-values, new confirmatory endpoint, or retrospective revision of the
lexical and causal decisions. The current title and integration are provisional.

\subsection{Continuation wording and numerical consistency}

The first design crosses two goals with two cues, separately for emoji and
previously recorded word alternatives. It scores scenario-specific
continuation pairs and two later formulations: ``This meets my goal.'' versus
``This does not meet my goal.'', and ``This is suitable for my goal.'' versus
``This is unsuitable for my goal.'' Leading spaces are part of the scored
continuations. All use Llama 3.2 3B base, BF16 Transformers/CUDA, eager
attention, without a chat template or model-weight training. The pinned
\texttt{mlx-community/Llama-3.2-3B-bf16} artifact is loaded through Transformers;
its identifier does not imply MLX execution. Its revision is
\texttt{60a99aaf43164077157d64bf909b7b61143c6a6d}.

An earlier execution version produced unequal scores for identical common
continuation prefixes in 355/576 cue-bearing pairs and 17/48 absent-cue pairs.
The matched-execution revision pads candidate sequences to equal length,
normalizes the same FP32 logit rectangle, and gathers valid continuation
positions. It yields exact common-prefix equality for all 624 pairs; 1,248
scientific sequence scores and 64 historical checks comprise 1,312 forwards.
The results below use this complete revised inventory. Earlier records remain
available. The check qualifies these numerical comparisons without identifying
a general BF16 cause or establishing semantic validity.

Let $s(c,v)=\log P(A\mid c,v)-\log P(B\mid c,v)$ be the full-continuation
score difference. With the recorded author orientation $\eta\in\{-1,+1\}$,
the interaction is
\[
H=\eta\{[s(p,a)-s(p,b)]-[s(q,a)-s(q,b)]\}.
\]
Coding $C,V\in\{-1,+1\}$, the four oriented scores have the exact
representation $\eta s=b_0+b_C C+b_V V+b_{CV}CV$, with $H=4b_{CV}$.
Positive interaction requires $b_{CV}>0$; the expected cue-effect crossover
requires $b_{CV}>|b_V|$. Author alignment in every input requires
$\min_{C,V} CV\,\eta s(C,V)>0$. All-four alignment implies the crossover,
which implies positive interaction. This is an accounting of the four
observed scores, not a decomposition of neural causes.

Table~\ref{tab:development-formulations} retains all formulations and both
cue types. Positive interaction across every formulation occurs on 22/24
emoji and 23/24 word scenarios; crossover across every formulation occurs on
8/24 and 14/24. All-four author alignment across every formulation occurs on
0/24 for each. Candidate polarity changes on 55/96 emoji and 56/96 word inputs.
The alternatives jointly vary vocabulary, specificity, morphology, and length.
Neither their semantic equivalence nor the authored labels has independent
human validation. Absent-cue subtraction changes offsets and context main
effects while leaving $H$ and cue-effect crossover unchanged algebraically;
it defines a different estimand rather than corrected accuracy.

\subsection{Three questions and an author-only tryout}

After an earlier author tryout, a revised design separates three propositions:
whether an accurately reported property matches a requested property;
whether a naturalistic message suggests a stated property or affect; and
whether that message alone establishes a specified wider outcome. The affect
questions concern visible expression in the bounded task and affect in the
naturalistic task. The changes therefore extend beyond instruction wording.

One previously exposed author saw 24 scenes, one condition per scenario:
12 bounded property questions and 12 naturalistic scenes with two questions
each, giving 36 substantive judgments. Emoji and word cues each appear in
12 scenes. Assignment used a source-only seed and did not select by model
agreement. Both naturalistic questions are visible together, allowing priming.
A separate assessability question precedes each Yes/No judgment. An undecidable
response leaves the substantive answer empty; it is never recoded to No.
English scenes accompany bilingual interface guidance. A human-only reading
was recorded before the current assignment/model comparison; known prior
project exposure prevents an independent-blinding claim.

% Generated by build_development_summary_20261010.py; do not edit.
\begin{table}[H]
\centering\small
\begin{tabular}{llrrrr}
\toprule
Formulation & Cue & Mean $H$ & $H>0$ & Crossover & All four \\
\midrule
Scenario-specific & Emoji & 1.3286 & 23/24 & 13/24 & 0/24 \\
Scenario-specific & Word & 2.2694 & 23/24 & 18/24 & 3/24 \\
Explicit negation & Emoji & 1.2329 & 23/24 & 12/24 & 1/24 \\
Explicit negation & Word & 1.6349 & 23/24 & 18/24 & 4/24 \\
Adjective pair & Emoji & 1.3695 & 24/24 & 15/24 & 2/24 \\
Adjective pair & Word & 2.2984 & 24/24 & 21/24 & 5/24 \\
\bottomrule
\end{tabular}
\caption{Matched-execution response scoring on the same 24 development scenarios. $H$ is an author-oriented context--cue interaction in nats. Crossover requires the expected reversal of the cue effect. All four requires author-aligned choices in every crossed input; authored labels are not independent human judgments.}
\label{tab:development-formulations}
\end{table}
\begin{table}[H]
\centering\small
\begin{tabular}{lrrrrr}
\toprule
Question & Yes & No & Undecidable & Binary & Agreement \\
\midrule
Property match & 7 & 5 & 0 & 12/12 & 7/12 \\
Message interpretation & 3 & 6 & 3 & 9/12 & 6/9 \\
Outcome evidence & 2 & 3 & 7 & 5/12 & 3/5 \\
\bottomrule
\end{tabular}
\caption{One previously exposed author's development tryout: 24 scenes, 36 judgments, no missing answers or abstentions. Agreement uses the exact positive-worded question and binary human judgments only. It is response-selected agreement, not accuracy; undecidable responses remain separate.}
\label{tab:development-author}
\end{table}


All scenes were completed, with 26 binary and ten undecidable judgments.
Table~\ref{tab:development-author} reports only aggregates. The model
comparison uses the exact positive-worded question and binary human answers,
giving 16/26 agreement on a response-selected subset. Model scoring forces
Yes/No while the human interface permits undecidability. Every one of the
36 selected model choices is also reproduced by a post-hoc constant rule:
Yes for property match and No for both naturalistic questions. On 32/36
selected questions, reversing true/false wording changes the implied truth.
The tryout locates question-design difficulties; it supplies no independent
human accuracy, population estimate, or causal improvement over the earlier
interface. Individual answers, assignments, and notes remain private outside
the repository, shared storage, and compute worker.

\subsection{Full model inventory and explicit controls}

Appendix~\ref{app:contextual-inputs} gives complete source-defined model inputs,
including the crossed cues, word alternatives, and accurate-completion controls.

The revised acquisition was specified before its outputs. For each of the
24 scenarios it crosses two contexts, five cue conditions (two emoji, two
words, one absent), three propositions, and two question polarities, giving
1,440 candidate pairs. Explicit accurate-completion and unknown-outcome
reports add 192 pairs. Thus 1,632 pairs require 3,264 scored sequences;
32 historical repeat checks bring the total to 3,296 forwards. All planned
sequences complete, every shared prefix is exact, and every historical check
matches. The recorded RTX 5090 acquisition takes 638.745 seconds, with peak
Torch reserved memory of 6,318 MiB. These counts do not add new scenario N.

Questions begin ``Is it true that'' or ``Is it false that'' with the same
following proposition. Full answers are \texttt{The answer is yes.}
and \texttt{The answer is no.}, each with an initial space and five tokens.
The literal score is $s=\log P(\mathrm{Yes\ answer})-
\log P(\mathrm{No\ answer})$; positive $s$ selects Yes, negative $s$ selects No,
and zero remains a tie. Negating the question reverses the answer needed to
preserve the underlying proposition. Normalized scores therefore use $-s$
for negative questions. No human undecidability option is modeled.

\begin{table}[H]
\centering\small
\begin{tabular}{lrrr}
\toprule
Positive-worded question & Emoji & Word & Absent cue \\
\midrule
Property match: Yes & 94/96 & 95/96 & 48/48 \\
Message interpretation: No & 94/96 & 90/96 & 46/48 \\
Outcome evidence: No & 95/96 & 93/96 & 48/48 \\
\bottomrule
\end{tabular}
\caption{Full-inventory answer preferences. Each 96 comprises 24 scenarios
by two contexts by two cues. Absent-cue conditions are shared, not duplicated
as two cue alternatives. These are preferences, not correctness counts.}
\label{tab:development-bias}
\end{table}

Across 816 paired true/false questions, 106 imply consistent truth and 710
imply inconsistent truth, with zero primary-score ties. All 48
positive-worded accurate-completion controls prefer No; their negative
versions also prefer No in 48/48 cases. Unknown-outcome controls prefer Yes
in 2/48 positive and 5/48 negative versions. Agreement on one question
polarity alone can therefore conceal contradictory implied interpretations.
The observations limit this checkpoint/prompt/readout combination. They do
not establish model-wide inability to understand negation or its internal cause.

\begin{table}[H]
\centering\small
\begin{tabular}{lrrrr}
\toprule
& \multicolumn{2}{c}{Emoji} & \multicolumn{2}{c}{Word} \\
Question & Positive & Negative$^\dagger$ & Positive & Negative$^\dagger$ \\
\midrule
Property match & 0.405 & 0.064 & 0.596 & 0.119 \\
Message interpretation & 0.392 & $-0.086$ & 0.829 & $-0.224$ \\
Outcome evidence & 0.059 & $-0.010$ & 0.123 & $-0.031$ \\
\bottomrule
\end{tabular}
\caption{Mean context--cue contrasts across 24 scenarios, in nats:
$[s(p,a)-s(p,b)]-[s(q,a)-s(q,b)]$.
$^\dagger$Negative-worded questions are sign-normalized to the same truth
direction. These descriptive interactions do not establish semantic accuracy
or an emoji advantage.}
\label{tab:development-contrasts}
\end{table}

For message interpretation, positive-worded contrasts are positive in 23/24
emoji and 24/24 word scenarios. Both become 7/24 after negative-wording
normalization. Continuous context sensitivity thus coexists with substantial
question dependence and failed explicit controls.

\paragraph{Secondary answer-token diagnostic.}
The shared three-token answer prefix cancels exactly in all 1,632 pairs.
The remaining sequence margin equals the yes/no-token difference plus the
following period difference. The period contributes a mean $-0.232795$ nats.
Using only the answer token reverses 103 non-tied signs and introduces
191 ties. Paired question polarities then comprise 78 consistent, 553
inconsistent, and 185 pairs with at least one tie. All 48 positive-worded
completion controls still prefer No. This post-result diagnostic leaves the
full-sequence primary estimator unchanged; omitting ties or replacing the
estimator would not qualify the present measurement.

\subsection{Complete source-defined contextual inputs}
\label{app:contextual-inputs}

These examples come from the public-source scenario \texttt{thermal\_01}
and its source-defined model inventory, not from a participant packet,
assignment, or response. The two contexts change Mira's stated goal; the
two cues are fire and ice. Each displayed \texttt{U+1F525} stands for the
single fire character and each \texttt{U+1F9CA} for the single ice character.
The model receives those Unicode characters, never their codepoint spellings.
Line wrapping below is typographic; the breaks before \texttt{Question:}
and \texttt{Answer:} are actual input newlines. The word comparators replace
only the character with the recorded alternative \texttt{hot} or
\texttt{cold}, respectively. The absent-cue condition uses the literal
\texttt{[no cue provided]}.

\newenvironment{contextinput}{\par\smallskip\noindent
  \begin{minipage}{\linewidth}\small\ttfamily\raggedright
  \hrule\smallskip}{\par\smallskip\hrule\end{minipage}\par\smallskip}

\paragraph{Complete crossed message-interpretation inputs.}
All four positive-worded prompts are shown. The corresponding negative
question replaces only \texttt{Is it true that} with
\texttt{Is it false that}. Both cue types are evaluated in both contexts;
their author-intended temperature interpretation has not been established
by independent human judgments.

% exact-input: thermal_01::p::emoji::a::message_reading::positive
\begin{contextinput}
Mira has come in from the snow and wants a drink that will warm her up. The cafe sends a message about the drink: U+1F525\\
Question: Is it true that the message suggests that the drink is hot?\\
Answer:
\end{contextinput}

% exact-input: thermal_01::p::emoji::b::message_reading::positive
\begin{contextinput}
Mira has come in from the snow and wants a drink that will warm her up. The cafe sends a message about the drink: U+1F9CA\\
Question: Is it true that the message suggests that the drink is hot?\\
Answer:
\end{contextinput}

% exact-input: thermal_01::q::emoji::a::message_reading::positive
\begin{contextinput}
Mira has come in from the heat and wants a drink that will cool her down. The cafe sends a message about the drink: U+1F525\\
Question: Is it true that the message suggests that the drink is cold?\\
Answer:
\end{contextinput}

% exact-input: thermal_01::q::emoji::b::message_reading::positive
\begin{contextinput}
Mira has come in from the heat and wants a drink that will cool her down. The cafe sends a message about the drink: U+1F9CA\\
Question: Is it true that the message suggests that the drink is cold?\\
Answer:
\end{contextinput}

The same naturalistic scenes also use the outcome question
\texttt{Is it true that the message alone establishes that Mira has warmed up?}
in the snow context and
\texttt{Is it true that the message alone establishes that Mira has cooled down?}
in the heat context, each with its corresponding \texttt{false} wording.
This changes the measured proposition from a reported drink property to
evidence about Mira's resulting condition.

\paragraph{Bounded property-matching input.}
The bounded arm adds a stipulated-accuracy instruction. This complete
snow/fire input illustrates the distinction from the naturalistic arm.
For the cold request, \texttt{hot} on the first line becomes \texttt{cold};
the reported cue still varies independently.

% exact-input: thermal_01::p::emoji::a::local_match::positive
\begin{contextinput}
The requested drink temperature is: hot.\\
Report of the drink temperature: U+1F525\\
For this task, assume the report accurately describes the drink temperature. It describes a state or category, not an instruction or a button action. Judge only the requested drink temperature; other goals and practical outcomes are outside this question. If the cue's meaning is unclear, the report's assumed accuracy does not resolve that meaning.\\
Question: Is it true that the reported drink temperature matches the requested drink temperature?\\
Answer:
\end{contextinput}

\paragraph{Accurate-completion controls and scored answers.}
These two prompts share an explicit, stipulated-accurate outcome report.
The intended answers are Yes and No, respectively. They test whether this
prompt and scoring arrangement distinguishes the reversed questions without
depending on an emoji's interpretation. Their recorded failure is scoped
to that arrangement.

% exact-input: thermal_01::p::complete_report::outcome_evidence::positive
\begin{contextinput}
For this question, assume this report is accurate: Mira has warmed up.\\
Question: Is it true that the message alone establishes that Mira has warmed up?\\
Answer:
\end{contextinput}

% exact-input: thermal_01::p::complete_report::outcome_evidence::negative
\begin{contextinput}
For this question, assume this report is accurate: Mira has warmed up.\\
Question: Is it false that the message alone establishes that Mira has warmed up?\\
Answer:
\end{contextinput}

Every prompt above is scored with the two exact continuations
\verb| The answer is yes.| and \verb| The answer is no.|.
The initial ASCII space and final period are included in the score. No
continuation text is generated and no chat template is inserted.
The unknown-outcome anchor instead uses
\texttt{The message explicitly says: I do not know whether the practical outcome happened.}
with the same two questions and continuations.

The scenario source is \path{data/contextual_emoji_development_v0.jsonl};
\path{scripts/contextual_emoji_two_track_v1.py} supplies the model
input construction. A regression check reconstructs all seven displayed
inputs, both word substitutions in both contexts, question reversals,
and the scored continuations directly from those public sources. It neither
opens private author records nor recomputes a model score.


\subsection{Qualification and remaining questions}
\label{app:contextual-next}

The extension motivated two qualification steps. For people, short questions
and explicit practice cases should separate message interpretation, event
truth, and evidential sufficiency while retaining uncertainty. For models,
simple text-only complete, incomplete, and unknown cases should test direct
versus nested questions, proposition reversals, answer-label mappings, and
constant-response comparators before evaluating an emoji effect.
Appendix~\ref{app:qualification} reports that text-only follow-up on new,
source-authored facts, using scoring and pass criteria specified before
execution. Model changes and independent human evaluation remain separate steps.

A stable cue meaning with a changing goal can explain a context--cue
interaction. Distinguishing that case from a changed interpretation of the
same cue requires new scenario/template families and independent evaluators.
Only a qualified output measure can support the intended next natural-state
causal study. No new sample size, threshold, recruitment, or intervention is
committed for a contextual-interpretation or causal follow-up by this revision.

The aggregate export rechecks saved scores, human/model joins, token diagnostics,
and their digests. Independent replay of human counts requires the private
records, which remain outside the September archive. These checks establish
faithful accounting. Semantic validation remains open. The October 10 aggregate
export added no model forwards, generation, or inferential tests; the later
acquisition in Appendix~\ref{app:qualification} is recorded separately.
